% Fichier complété par tools/complete_missing_corrections.py.
% Source : NUM/NUM-03/1023_Euler/1023_Euler.tex
% Statut : rédigé (1 question(s)).
\begin{corrigebox}[Corrigé rédigé]
\CorrigeQuestion{1}
Pour une équation $\dot y=f(t,y)$ et un pas $h$, le schéma d'Euler
        explicite est $t_{n+1}=t_n+h$ et
        \[\boxed{y_{n+1}=y_n+h f(t_n,y_n)}.\]
        Pour un système d'ordre supérieur, on introduit les variables d'état
        $x_1=y$, $x_2=\dot y$, etc., puis on applique cette formule à chaque composante.
        Exemple pour $\ddot y=g(t,y,\dot y)$ :
        $v_{n+1}=v_n+h g(t_n,y_n,v_n)$ et $y_{n+1}=y_n+h v_n$.
        Les conditions initiales fixent $y_0$ et $v_0$; le pas est diminué jusqu'à
        convergence des résultats.
\end{corrigebox}
